\chapter{Stopping Times}

In many financial applications (e.g., American options, optimal stopping problems), actions depend on random events. A fundamental concept is the \textbf{stopping time}, which models a decision to stop or take action based \textit{only} on currently available information.

\section{Definition and Intuition}

A random variable $\tau : \Omega \to I \cup \{+\infty\}$ is a time index that is random. However, for it to be useful in decision making, observing the process up to time $t$ must be sufficient to determine if $\tau$ has occurred by time $t$.

\begin{definition}[Stopping Time]
A random variable $\tau$ taking values in $I \cup \{+\infty\}$ is a \textbf{stopping time} with respect to the filtration $\mathbb{F} = (\mathcal{F}_t)_{t \in I}$ if for every $t \in I$:
\[
\{\tau \le t\} \in \mathcal{F}_t.
\]
\end{definition}

\textbf{Intuition}: The event $\{\tau \le t\}$ means "the stopping time has occurred by time $t$". The definition requires this event to be observable given the information $\mathcal{F}_t$.
If $\tau$ represents the time to sell a stock, the decision to sell \textit{at or before time $t$} must rely only on stock prices up to time $t$, not on future prices.

\begin{proposition}
If $I = \mathbb{N}$ (discrete time), then $\tau$ is a stopping time if and only if:
\[
\{\tau = t\} \in \mathcal{F}_t, \quad \forall t \in I.
\]
\end{proposition}
\begin{proof}
$(\Rightarrow)$ If $\tau$ is a stopping time, then $\{\tau \le t\} \in \mathcal{F}_t$ and $\{\tau \le t-1\} \in \mathcal{F}_{t-1} \subseteq \mathcal{F}_t$.
Since $\{\tau = t\} = \{\tau \le t\} \setminus \{\tau \le t-1\}$, it is the difference of two sets in $\mathcal{F}_t$, hence in $\mathcal{F}_t$.
$(\Leftarrow)$ If $\{\tau = k\} \in \mathcal{F}_k$ for all $k$, then
\[
\{\tau \le t\} = \bigcup_{k=0}^t \{\tau = k\}.
\]
Since $\{\tau = k\} \in \mathcal{F}_k \subseteq \mathcal{F}_t$ for $k \le t$, the finite union is in $\mathcal{F}_t$.
\end{proof}

\section{Examples of Stopping Times}

\subsection{First Hitting Time}
Let $A$ be a Borel set in the state space $E$ (e.g., $A = [K, \infty)$ for a price threshold).
Basically, the \textbf{first hitting time} of $A$ is defined as:
\[
H_A(\omega) = \inf \{ t \in I : X_t(\omega) \in A \}.
\]
(Convention: $\inf \emptyset = +\infty$).

\begin{proposition}
If $X$ is an adapted process, then $H_A$ is a stopping time.
\end{proposition}
\begin{proof}
We check the condition $\{\tau \le t\} \in \mathcal{F}_t$:
\begin{align*}
\{H_A \le t\} &= \{ \exists k \in \{0, \dots, t\} \text{ s.t. } X_k \in A \} \\
&= \bigcup_{k=0}^t \{ X_k \in A \}.
\end{align*}
Since $X$ is adapted, $\{ X_k \in A \} \in \mathcal{F}_k \subseteq \mathcal{F}_t$ for all $k \le t$.
Since $\mathcal{F}_t$ is a $\sigma$-algebra, it is closed under finite unions. Thus, $\{H_A \le t\} \in \mathcal{F}_t$.
\end{proof}

\subsection{Counter-Example: Last Hitting Time}
Define $L_A = \sup \{ t \in I : X_t \in A \}$. This is generally \textbf{not} a stopping time.
To know if $L_A \le t$, we need to know that the process will \textit{never} visit $A$ again after time $t$. This requires future information, which is not available in $\mathcal{F}_t$.

\section{Stopped Processes}

Given a process $X$ and a stopping time $\tau$, we often want to consider the process "frozen" at time $\tau$. This is crucial for analyzing strategies that exit a market or stop a game.

\begin{definition}[Stopped Process]
The \textbf{stopped process} $X^\tau = (X^\tau_t)_{t \in I}$ is defined by:
\[
X^\tau_t(\omega) = X_{t \wedge \tau(\omega)}(\omega) = \begin{cases}
X_t(\omega) & \text{if } t < \tau(\omega) \\
X_{\tau(\omega)}(\omega) & \text{if } t \ge \tau(\omega)
\end{cases}
\]
where $a \wedge b = \min(a, b)$.
\end{definition}

\begin{proposition}
If $X$ is adapted and $\tau$ is a stopping time, then the stopped process $X^\tau$ is adapted.
\end{proposition}
\begin{proof}
We need to show that $X^\tau_t$ is $\mathcal{F}_t$-measurable for each $t$.
We can write:
\[
X^\tau_t = X_t \cdot \mathbb{I}_{\{t < \tau\}} + \sum_{k=0}^t X_k \cdot \mathbb{I}_{\{\tau = k\}}.
\]
Let's analyze each term:
1. $\{t < \tau\} = \{\tau \le t\}^c \in \mathcal{F}_t$ (since $\tau$ is a stopping time). $X_t$ is $\mathcal{F}_t$-measurable. Thus the product is $\mathcal{F}_t$-measurable.
2. For each $k \le t$, regarding the term in the sum:
   $\{\tau = k\} \in \mathcal{F}_k \subseteq \mathcal{F}_t$.
   $X_k$ is $\mathcal{F}_k$-measurable $\implies \mathcal{F}_t$-measurable.
   Thus $X_k \cdot \mathbb{I}_{\{\tau = k\}}$ is $\mathcal{F}_t$-measurable.
Sum of measurable functions is measurable.
Therefore, $X^\tau_t$ is $\mathcal{F}_t$-measurable.
\end{proof}
